% GENERATED FILE. Edit canonical lessons, then run npm run generate:books.
% canonical-source-hash: fnv1a64:8b733c26a5507523
% canonical-lessons: TE-C123-puli, TE-C123-simham, TE-C123-elugubanti, TE-C123-koti, TE-C123-kundelu

\chapter{Tiger, Lion, Bear, Monkey, Rabbit}
\label{ch:te-tiger-lion-bear-monkey-rabbit}
\hlchaptermodality{123}

\begin{chapteropening}
\textbf{By the end of this chapter:} \emph{I can say a tiger, a lion, a bear, a monkey and a rabbit.}

Complete the last lesson of chapter 123: I can say a tiger, a lion, a bear, a monkey and a rabbit.
\end{chapteropening}

\section[\texorpdfstring{\te{పులి} (puli)}{పులి (puli)}]{\te{పులి} (puli) — a tiger}

\label{lesson:TE-C123-puli}



\begin{quote}
Before the new one: say the Telugu for a horse, then the Telugu for an elephant.
\end{quote}

\subsection*{You'll want to know: \te{పులి}}
\textbf{\te{పులి}} — \emph{puli} — \textquotedblleft{}a tiger\textquotedblright{}.

\begin{cousinweb}
\textbf{\te{పులి}} (\emph{puli}) is native Dravidian, one word with Tamil \emph{puli} and Kannada \emph{puli} (DEDR 4307). DEDR compares Prakrit \emph{pulli}, \textquotedblleft{}tiger\textquotedblright{}.
\end{cousinweb}

\begin{grammarlens}[title={What you've built}]
One more word for this chapter.
\end{grammarlens}

\subsection*{Your turn}
\begin{itemize}
\raggedright
  \item \emph{Say it:} \emph{puli}
  \item \emph{Say it:} \emph{puli}, once more
  \item \emph{Say it:} say \emph{puli}
  \item \emph{Recall:} say the Telugu for a horse, then the Telugu for an elephant, then say \emph{puli} again
\end{itemize}

\begin{tcolorbox}[breakable,colback=teal!4,colframe=teal!35!black,title={Before you move on}]
What does \te{పులి} mean? (\textquotedblleft{}A tiger\textquotedblright{}.) Say it once more.
\end{tcolorbox}
\section[\texorpdfstring{\te{సింహం} (siṁhaṁ)}{సింహం (siṁhaṁ)}]{\te{సింహం} (siṁhaṁ) — a lion}

\label{lesson:TE-C123-simham}



\begin{quote}
Before the new one: say the Telugu for an elephant, then the Telugu for a tiger.
\end{quote}

\subsection*{You'll want to know: \te{సింహం}}
\textbf{\te{సింహం}} — \emph{siṁhaṁ} — \textquotedblleft{}a lion\textquotedblright{}.

\begin{cousinweb}
\textbf{\te{సింహం}} (\emph{siṁhaṁ}) is Sanskrit \emph{siṁha}, \textquotedblleft{}lion\textquotedblright{}, taken over whole as a learned loan, a tatsama, with the Telugu noun ending \emph{-ṁ} (CDIAL 13384).
\end{cousinweb}

\begin{grammarlens}[title={What you've built}]
One more word for this chapter.
\end{grammarlens}

\subsection*{Your turn}
\begin{itemize}
\raggedright
  \item \emph{Say it:} \emph{siṁhaṁ}
  \item \emph{Say it:} \emph{siṁhaṁ}, once more
  \item \emph{Say it:} say \emph{siṁhaṁ}
  \item \emph{Recall:} say the Telugu for an elephant, then the Telugu for a tiger, then say \emph{siṁhaṁ} again
\end{itemize}

\begin{tcolorbox}[breakable,colback=teal!4,colframe=teal!35!black,title={Before you move on}]
What does \te{సింహం} mean? (\textquotedblleft{}A lion\textquotedblright{}.) Say it once more.
\end{tcolorbox}
\section[\texorpdfstring{\te{ఎలుగుబంటి} (elugubaṇṭi)}{ఎలుగుబంటి (elugubaṇṭi)}]{\te{ఎలుగుబంటి} (elugubaṇṭi) — a bear}

\label{lesson:TE-C123-elugubanti}



\begin{quote}
Before the new one: say the Telugu for a tiger, then the Telugu for a lion.
\end{quote}

\subsection*{You'll want to know: \te{ఎలుగుబంటి}}
\textbf{\te{ఎలుగుబంటి}} — \emph{elugubaṇṭi} — \textquotedblleft{}a bear\textquotedblright{}.

\begin{cousinweb}
\textbf{\te{ఎలుగుబంటి}} (\emph{elugubaṇṭi}) is built on native \emph{elugu}, \textquotedblleft{}bear\textquotedblright{}, one word with Tamil \emph{eṇku}, \textquotedblleft{}bear\textquotedblright{} (DEDR 857).
\end{cousinweb}

\begin{grammarlens}[title={What you've built}]
One more word for this chapter.
\end{grammarlens}

\subsection*{Your turn}
\begin{itemize}
\raggedright
  \item \emph{Say it:} \emph{elugubaṇṭi}
  \item \emph{Say it:} \emph{elugubaṇṭi}, once more
  \item \emph{Say it:} say \emph{elugubaṇṭi}
  \item \emph{Recall:} say the Telugu for a tiger, then the Telugu for a lion, then say \emph{elugubaṇṭi} again
\end{itemize}

\begin{tcolorbox}[breakable,colback=teal!4,colframe=teal!35!black,title={Before you move on}]
What does \te{ఎలుగుబంటి} mean? (\textquotedblleft{}A bear\textquotedblright{}.) Say it once more.
\end{tcolorbox}
\section[\texorpdfstring{\te{కోతి} (kōti)}{కోతి (kōti)}]{\te{కోతి} (kōti) — a monkey}

\label{lesson:TE-C123-koti}



\begin{quote}
Before the new one: say the Telugu for a lion, then the Telugu for a bear.
\end{quote}

\subsection*{You'll want to know: \te{కోతి}}
\textbf{\te{కోతి}} — \emph{kōti} — \textquotedblleft{}a monkey\textquotedblright{}.

\begin{cousinweb}
\textbf{\te{కోతి}} (\emph{kōti}) is native Dravidian, one word with Tamil \emph{kuraṅku} and Kannada \emph{kōti} (DEDR 1769).
\end{cousinweb}

\begin{grammarlens}[title={What you've built}]
One more word for this chapter.
\end{grammarlens}

\subsection*{Your turn}
\begin{itemize}
\raggedright
  \item \emph{Say it:} \emph{kōti}
  \item \emph{Say it:} \emph{kōti}, once more
  \item \emph{Say it:} say \emph{kōti}
  \item \emph{Recall:} say the Telugu for a lion, then the Telugu for a bear, then say \emph{kōti} again
\end{itemize}

\begin{tcolorbox}[breakable,colback=teal!4,colframe=teal!35!black,title={Before you move on}]
What does \te{కోతి} mean? (\textquotedblleft{}A monkey\textquotedblright{}.) Say it once more.
\end{tcolorbox}
\section[\texorpdfstring{\te{కుందేలు} (kundēlu)}{కుందేలు (kundēlu)}]{\te{కుందేలు} (kundēlu) — a rabbit}

\label{lesson:TE-C123-kundelu}



\begin{quote}
Before the new one: say the Telugu for a bear, then the Telugu for a monkey.
\end{quote}

\subsection*{You'll want to know: \te{కుందేలు}}
\textbf{\te{కుందేలు}} — \emph{kundēlu} — \textquotedblleft{}a rabbit\textquotedblright{}.

\begin{cousinweb}
\textbf{\te{కుందేలు}} (\emph{kundēlu}) is native Dravidian. DEDR 1728 lists the Telugu word for \textquotedblleft{}hare\textquotedblright{} under the verb for squatting, with Tamil \emph{kuntu}, \textquotedblleft{}to sit on the heels\textquotedblright{}.
\end{cousinweb}

\begin{grammarlens}[title={What you've built}]
One more word for this chapter.
\end{grammarlens}

\subsection*{Your turn}
\begin{itemize}
\raggedright
  \item \emph{Say it:} \emph{kundēlu}
  \item \emph{Say it:} \emph{kundēlu}, once more
  \item \emph{Say it:} say \emph{kundēlu}
  \item \emph{Recall:} say the Telugu for a bear, then the Telugu for a monkey, then say \emph{kundēlu} again
\end{itemize}

\begin{tcolorbox}[breakable,colback=teal!4,colframe=teal!35!black,title={Before you move on}]
What does \te{కుందేలు} mean? (\textquotedblleft{}A rabbit\textquotedblright{}.) Say it once more.
\end{tcolorbox}
